%% 5-16 積分の道具箱：置換積分と部分積分（第11回：ページ制限なし。積分の道具を1段ずつ書き足した）
\newcommand\bkvpmidashi[1]{\par\Needspace{6\baselineskip}\noindent\textbf{#1}\par}
\bkchapter{5}{16}{積分の道具箱：\\置換積分と部分積分}{\bklevel{4}}{積分の道具箱：置換積分と部分積分}
\bkmichishirube{5-6,5-11,5-14}{6-4,9-1,9-5}

\begin{goal}
\begin{itemize}[label={\color{bkTheme}●}, leftmargin=1.6\zw, itemsep=0pt, topsep=1pt]
\item \bklevel{4}\enspace 合成関数の微分を逆に使って、置換積分ができる。
\item \bklevel{4}\enspace 積の微分を逆に使って、部分積分ができる。
\item \bklevel{4}\enspace 流れた電気の量や円の面積など、表にない形の積分を求められる。
\end{itemize}
\end{goal}

\begin{hatena}
\noindent\begin{minipage}[c]{0.535\linewidth}
ロボットのモーターのスイッチを入れると、電流が $i=2t\,e^{-t}$（A、$t$ は秒）のように、ふえてからへっていきました。はじめの3秒間に流れた電気の量（グラフの面積）は、何C（クーロン）でしょう。

流れた電気の量は、電流の積分です（5-12）。でも、$t$ と $e^{-t}$ のかけ算の積分は、まだ表にありません。
\end{minipage}\hfill
\begin{minipage}[c]{0.425\linewidth}
\begin{tikzpicture}
\begin{axis}[bkplot, width=56mm, height=38mm, xmin=0, xmax=4.2, ymin=0, ymax=0.95, xtick={1,2,3,4}, ytick={0.5}, xlabel={$t$}, ylabel={$i$}]
  \addplot[bkcurve, domain=0:3, samples=60, fill=bkYellow!45, draw=none] {2*x*exp(-x)} \closedcycle;
  \addplot[bkcurve, domain=0:4.1, samples=80] {2*x*exp(-x)};
  \node[font=\sffamily\bfseries\small] at (axis cs:1.5,0.35) {？};
\end{axis}
\end{tikzpicture}
\end{minipage}
\end{hatena}

積分は「微分すると、その式になる関数」をさがす計算でした（5-11）。この章では、微分の2つの道具（5-6）を「逆向き」に使って、さがし方の道具を2つ作ります。

\clearpage
%% ---------------- やってみよう・思い出そう ----------------
\begin{yattemiyou}[表計算で、面積をおおよそ求めよう]
5-10と同じように、0から3を、はば0.5の細い長方形6本に分け、長方形の面積をたします。高さは、それぞれの区間のまん中の $t$ での $i$ の値にします。
\begin{bkcenter}
\bksheet{4}{22}{%
1 & 区間のまん中 $t$ & $i=2t\,e^{-t}$ & はば & 長方形の面積 \\
2 & 0.25 & 0.3894 & 0.5 & 0.1947 \\
3 & 0.75 & 0.7085 & 0.5 & 0.3543 \\
4 & 1.25 & 0.7163 & 0.5 & 0.3581 \\
5 & 1.75 & 0.6082 & 0.5 & 0.3041 \\
6 & 2.25 & 0.4743 & 0.5 & 0.2371 \\
7 & 2.75 & 0.3516 & 0.5 & 0.1758 \\
8 & & & 合計 & 1.6242 \\}
\bktabcaption{はばを0.1にすると1.6026、0.01にすると1.6017。1.60…に近づいていく。}
\end{bkcenter}
\end{yattemiyou}

\noindent 表計算なら、およそ1.60Cと分かります。でも、「1.60…」のもとになる、正確な式は何でしょう。この章の道具を使うと、答えは $2-8e^{-3}$ と分かります（⑩と例題8）。

\begin{omoidasou}
\bkrecall{5-4}{定数倍は前に出せる：$(cf)'=cf'$。積分でも $\int cf\,\dd x=c\int f\,\dd x$}
\bkrecall{5-6}{積の微分 $(fg)'=f'g+fg'$。合成関数の微分 $\dfrac{\dd y}{\dd x}=\dfrac{\dd y}{\dd u}\times\dfrac{\dd u}{\dd x}$（外側の微分×内側の微分）}
\bkrecall{5-8}{$x$ が少し（$\Delta x$）変わると、$y$ は約 $f'(x)\,\Delta x$ 変わる。$\Delta y\approx f'(x)\,\Delta x$}
\bkrecall{5-11}{微分して $f(x)$ になる関数 $F(x)$ を見つければ、$\displaystyle\int_a^bf(x)\,\dd x=F(b)-F(a)$。$\displaystyle\int x^n\,\dd x=\frac{x^{n+1}}{n+1}+C$}
\bkrecall{5-14}{$\displaystyle\int e^{ax}\,\dd x=\frac{1}{a}e^{ax}+C$、$\displaystyle\int\frac{1}{x}\,\dd x=\ln x+C$（$x>0$）}
\bkrecall{5-15}{$\displaystyle\int\sin t\,\dd t=-\cos t+C$、$\displaystyle\int\cos t\,\dd t=\sin t+C$、$\cos^2\theta=\dfrac{1+\cos2\theta}{2}$}
\end{omoidasou}

\Needspace{14\baselineskip}
%% ---------------- 数学の道具 ----------------
\bksection{doguu}{数学の道具}

\bksubsection{1}{積分は「微分の答え」から見つける}

\label{bk:5-16:hajime}5-11では、微分の表を逆に読んで、積分の表を作りました。
\begin{bkcenter}
\begin{bktable}{colspec={Q[c,m,wd=44mm] Q[c,m,wd=50mm]}, rowsep=2pt}
微分の式（左を微分すると右） & 積分の式（右を積分すると左） \\
$(x^3)'=3x^2$ & $\displaystyle\int3x^2\,\dd x=x^3+C$ \\
$(e^{2x})'=2e^{2x}$ & $\displaystyle\int2e^{2x}\,\dd x=e^{2x}+C$ \\
$(\sin x)'=\cos x$ & $\displaystyle\int\cos x\,\dd x=\sin x+C$ \\
\end{bktable}
\end{bkcenter}
積分したい式が、表の右の列にそのまま出ていれば、答えはすぐ分かります。では、$2x(x^2+1)^3$ や $x\,e^x$ はどうでしょう。こんな形は、表にはありません。

でも、5-6で、積の微分と合成関数の微分を使って、たくさんの式を微分しました。\textbf{その微分の答えを集めておけば、それを逆に読んで積分できる}はずです。この章の2つの道具は、この考え方から生まれます。
\begin{bkcenter}
\begin{bktable}{colspec={Q[l,m,wd=46mm] Q[l,m,wd=46mm]}, rowsep=2pt}
微分の道具（5-6） & 逆に使った積分の道具 \\
合成関数の微分 & 置換積分（②〜⑥） \\
積の微分 & 部分積分（⑦〜⑩） \\
\end{bktable}
\end{bkcenter}

\bksubsection{2}{合成関数の微分を、逆に読む}

\label{bk:5-16:gyaku}合成関数の微分で、$(x^2+1)^4$ を微分してみます。かっこの中を内側 $u=x^2+1$、4乗を外側 $u^4$ と見ます（5-6）。
\begin{align*}
\bigl\{(x^2+1)^4\bigr\}'&=\underbrace{4u^3}_{\text{外側の微分}}\times\underbrace{(x^2+1)'}_{\text{内側の微分}}=4(x^2+1)^3\times2x\\
&=8x(x^2+1)^3
\end{align*}
この式を逆に読むと、積分の式が1つ手に入ります。
\[ \int8x(x^2+1)^3\,\dd x=(x^2+1)^4+C \]
積分の中の $8x(x^2+1)^3$ を、色で分けて見てみましょう。

\begin{bkfigure}
\begin{tikzpicture}[x=1mm, y=1mm]
  \node[font=\large, inner sep=1pt, anchor=base east] (p1) at (0,0) {$\displaystyle\int\ 4\,$};
  \node[font=\large, inner sep=1pt, anchor=base west, text=bkBlue] (p2) at (p1.base east) {$(x^2+1)^3$};
  \node[font=\large, inner sep=1pt, anchor=base west] (p3) at (p2.base east) {$\,\times\,$};
  \node[font=\large, inner sep=1pt, anchor=base west, text=bkOrange!80!black] (p4) at (p3.base east) {$2x$};
  \node[font=\large, inner sep=1pt, anchor=base west] (p5) at (p4.base east) {$\,\dd x$};
  \node[bk node, draw=bkBlue, fill=bkBlue!6, font=\sffamily\footnotesize, align=center, anchor=north] (k1) at (-22,-14) {内側 $u=x^2+1$ が\\入った式};
  \node[bk node, draw=bkOrange, fill=bkOrange!6, font=\sffamily\footnotesize, align=center, anchor=north] (k2) at (40,-14) {内側の微分\\$(x^2+1)'=2x$};
  \draw[bk link, bkBlue] (k1.north) -- (p2.south);
  \draw[bk link, bkOrange] (k2.north) -- (p4.south);
\end{tikzpicture}
\bkfigcaption{積分の中に、「内側の式」と「内側の微分」が両方入っている。これが、合成関数の微分を逆に読める目じるし。}
\end{bkfigure}

図では、$8x$ を $4\times2x$ と分けて、「内側の微分 $2x$」が見えるようにしました。では、前の数が4でなく1のとき、$\displaystyle\int2x(x^2+1)^3\,\dd x$ はどうなるでしょう。$2x(x^2+1)^3$ は $8x(x^2+1)^3$ のちょうど $\frac{1}{4}$ です。だから答えも $\frac{1}{4}$ にすればよいはずです（定数倍は前に出せる、5-4）。
\[ \int2x(x^2+1)^3\,\dd x=\frac{1}{4}\int8x(x^2+1)^3\,\dd x=\frac{(x^2+1)^4}{4}+C \]
微分して確かめましょう。$\Bigl\{\dfrac{(x^2+1)^4}{4}\Bigr\}'=\dfrac{1}{4}\times8x(x^2+1)^3=2x(x^2+1)^3$。積分の中身にもどりました。

\bksubsection{3}{置換積分：$u$ におきかえる}

\label{bk:5-16:chikan}②では、答えを知っている式から逆にたどりました。こんどは、答えを知らないところから、手順だけで $\displaystyle\int2x(x^2+1)^3\,\dd x$ を求めます。
\begin{bkcenter}
% 第6回：3つの枠の大きさ（はば40mm）と上端をそろえ、中の字は上からそろえて組む。矢印は枠の名前どうしでつなぐ
\begin{tikzpicture}[x=1mm, y=1mm, st/.style={bk node, draw=bkTheme, fill=bkTheme!6, align=left, anchor=north west, font=\small}]
  \node[st] (a) at (0,0) {\parbox[t][23.5mm][t]{33mm}{\textbf{①おく}\par\vspace{2pt}内側を\par $u=x^2+1$ とおく。}};
  \node[st] (b) at (49.5,0) {\parbox[t][23.5mm][t]{33mm}{\textbf{②書きかえる}\par\vspace{2pt}$\dd u=2x\,\dd x$\par\vspace{2pt}$\displaystyle\int u^3\,\dd u$}};
  \node[st] (c) at (99,0) {\parbox[t][23.5mm][t]{33mm}{\textbf{③積分してもどす}\par\vspace{2pt}$\dfrac{u^4}{4}+C$\par\vspace{2pt}$=\dfrac{(x^2+1)^4}{4}+C$}};
  \draw[bk link, bkTheme] (a) -- (b);
  \draw[bk link, bkTheme] (b) -- (c);
\end{tikzpicture}
\end{bkcenter}
いちばん大事なのは、②の $\dd u=2x\,\dd x$ です。2通りの見方で説明します。

\bkvpmidashi{見方1：微分の記号を分数のように使う}
$u=x^2+1$ を $x$ で微分すると $\dfrac{\dd u}{\dd x}=2x$。両辺に $\dd x$ をかけると、$\dd u=2x\,\dd x$ になります。$\dfrac{\dd u}{\dd x}$ はほんとうの分数ではありませんが、もとは小さな変化の比 $\dfrac{\Delta u}{\Delta x}$ なので、このように計算してよいのです（5-6のコラム）。分数のように書きかえられるのは、ライプニッツ（5-3）が、計算しやすいように記号を工夫したおかげです。

\bkvpmidashi{見方2：数で確かめる}
$x$ が1から1.01まで、0.01だけ変わるとき、$u=x^2+1$ はどれだけ変わるでしょう。
\[ u(1.01)-u(1)=(1.01^2+1)-(1^2+1)=2.0201-2=0.0201 \]
一方、$2x\,\dd x$ に $x=1$、$\dd x=0.01$ を入れると、$2\times1\times0.01=0.02$。ほとんど同じです。「$x$ が少し（$\dd x$）変わると、$u$ はその $2x$ 倍だけ変わる」という意味です（5-8の $\Delta y\approx f'(x)\,\Delta x$ と同じ）。

\bkvpmidashi{書きかえ}
積分の中を、$u$ の式とそれ以外に分けて、1つずつ書きかえます。
\[ \int\underbrace{(x^2+1)^3}_{u^3}\ \underbrace{2x\,\dd x}_{\dd u}=\int u^3\,\dd u \]
$u$ だけの式になったので、5-11の「指数を1ふやして、その数でわる」で積分できます。$\dfrac{u^4}{4}+C$。最後に、$u$ を $x^2+1$ にもどして、$\dfrac{(x^2+1)^4}{4}+C$。②で求めた答えと同じになりました。

このように、積分の中の式を $u$ におきかえて、知っている形にする方法を\textbf{置換積分}といいます。「置換」は「おきかえる」という意味です。\bkidx{ちかんせきぶん}{置換積分}まとめると、
\[ \int(\text{$u$ の式})\times u'\,\dd x=\int(\text{$u$ の式})\,\dd u \]
です。$u'\,\dd x$（内側の微分×$\dd x$）が、まとめて $\dd u$ になります。

\bksubsection{4}{係数が合わないとき}

\label{bk:5-16:keisu}$\displaystyle\int x\,e^{x^2}\,\dd x$ を求めます。$e$ の肩の $x^2$ を内側と見て、$u=x^2$ とおきます。
\begin{enumerate}[label={\bkcircled[bkTheme]{\arabic*}}, leftmargin=2.4\zw, itemsep=1pt, topsep=2pt]
\item $u=x^2$ を微分すると $\dfrac{\dd u}{\dd x}=2x$ なので、$\dd u=2x\,\dd x$。
\item でも、積分の中には $x\,\dd x$ しかありません。$2x\,\dd x$ の半分です。そこで、$\dd u=2x\,\dd x$ の両辺を2でわって、$x\,\dd x=\dfrac{1}{2}\dd u$。
\item 書きかえると、$\displaystyle\int e^{u}\cdot\frac{1}{2}\,\dd u=\frac{1}{2}\int e^u\,\dd u=\frac{1}{2}e^u+C$（$\frac{1}{2}$ は定数なので前に出せる）。
\item $u$ をもどして、$\dfrac{1}{2}e^{x^2}+C$。
\end{enumerate}
確かめ：$\Bigl(\dfrac{1}{2}e^{x^2}\Bigr)'=\dfrac{1}{2}\times e^{x^2}\times2x=x\,e^{x^2}$。合っています。

\bkvpmidashi{使えないとき}
2や $\frac{1}{2}$ のような\textbf{数}のずれは、このように調整できます。でも、$\displaystyle\int e^{x^2}\,\dd x$ のように $x$ そのものが足りないときは、置換積分は使えません。$\dd x=\dfrac{\dd u}{2x}$ と書きかえても、分母に $x$ が残って、「$u$ だけの式」にならないからです（$x$ は定数ではないので、前に出せません）。

\bkvpmidashi{よく出る形：分母の微分が分子にある}
$\displaystyle\int\frac{2x}{x^2+1}\,\dd x$ では、分母 $x^2+1$ を微分すると分子の $2x$ になっています。そこで $u=x^2+1$（いつも正）とおくと、$\dd u=2x\,\dd x$ で、
\[ \int\frac{2x}{x^2+1}\,\dd x=\int\frac{1}{x^2+1}\cdot2x\,\dd x=\int\frac{1}{u}\,\dd u=\ln u+C=\ln(x^2+1)+C \]
（$\displaystyle\int\frac{1}{u}\,\dd u=\ln u+C$ は5-14。）確かめ：$\{\ln(x^2+1)\}'=\dfrac{1}{x^2+1}\times2x$（5-14と合成関数の微分）。

\bksubsection{5}{定積分の置換は、はんいも書きかえる}

\label{bk:5-16:hanni}定積分 $\displaystyle\int_0^12x(x^2+1)^3\,\dd x$ を求めます。

\bkvpmidashi{やり方1：$u$ のはんいで計算する}
$u=x^2+1$ とおくと、積分の中は $u$ だけの式になります。このとき、「$x$ が0から1まで」も、「$u$ がいくつからいくつまで」に書きかえます。積分の横軸が、$x$ から $u$ に変わったからです。
\begin{bkcenter}
\begin{bktable}{colspec={Q[l,m,wd=30mm] Q[c,m,wd=26mm] Q[c,m,wd=26mm]}, rowsep=2pt}
 & 下端 & 上端 \\
$x$ & 0 & 1 \\
$u=x^2+1$ & $0^2+1=1$ & $1^2+1=2$ \\
\end{bktable}
\end{bkcenter}
\[ \int_0^12x(x^2+1)^3\,\dd x=\int_1^2u^3\,\dd u=\Bigl[\,\frac{u^4}{4}\,\Bigr]_1^2=\frac{2^4}{4}-\frac{1^4}{4}=\frac{16}{4}-\frac{1}{4}=\frac{15}{4} \]
$x$ にもどさなくても、$u$ のはんいで計算すれば答えが出ます。

\bkvpmidashi{やり方2：$x$ にもどしてから計算する}
③の不定積分 $\dfrac{(x^2+1)^4}{4}$ を使って、$x$ のはんいで計算しても、
\[ \Bigl[\,\frac{(x^2+1)^4}{4}\,\Bigr]_0^1=\frac{2^4}{4}-\frac{1^4}{4}=\frac{15}{4} \]
で、同じ答えです。やり方1の $u=2$、$u=1$ は、$x=1$、$x=0$ のときの $x^2+1$ の値だったので、同じ計算になるのです。

\begin{bkfigure}
\begin{tikzpicture}
\begin{axis}[bkplot, name=L, width=60mm, height=38mm, xmin=0, xmax=1.15, ymin=0, ymax=19, xtick={0,1}, ytick={8,16}, xlabel={$x$}, title={\small $2x(x^2+1)^3$ を $x$：0〜1で}, title style={yshift=-1mm}]
  \addplot[bkcurve, domain=0:1, samples=50, fill=bkBlue!15] {2*x*(x^2+1)^3} \closedcycle;
  \addplot[bkcurve, domain=0:1.04, samples=50] {2*x*(x^2+1)^3};
\end{axis}
\begin{axis}[bkplot, at={(L.east)}, anchor=west, xshift=16mm, width=60mm, height=38mm, xmin=0.8, xmax=2.2, ymin=0, ymax=9.5, xtick={1,2}, ytick={4,8}, xlabel={$u$}, title={\small $u^3$ を $u$：1〜2で}, title style={yshift=-1mm}]
  \addplot[bkcurve2, domain=1:2, samples=50, fill=bkOrange!18] {x^3} \closedcycle;
  \addplot[bkcurve2, domain=0.8:2.1, samples=50] {x^3};
\end{axis}
\node[font=\sffamily\small] at ($(L.east)+(8mm,0)$) {$=$};
\end{tikzpicture}
\bkfigcaption{形も横のはんいもちがうのに、色の部分の面積はどちらも $\frac{15}{4}=3.75$（細い長方形1000本でたしても3.7500）。$x$ の目もりを $u=x^2+1$ の目もりにつけかえると、同じ面積を別の見方で計算できる。}
\end{bkfigure}
左の図で、$x$ のあいだ $\dd x$ の細い長方形（高さ $2x(x^2+1)^3$）の面積は、$(x^2+1)^3\times2x\,\dd x=u^3\,\dd u$ で、右の図の、はば $\dd u$ の細い長方形の面積と同じです。細い長方形が1本ずつ同じ面積なので、全部たした面積も同じになります。

\bksubsection{6}{$x$ を角でおきかえる：円の面積}

\label{bk:5-16:en}置換積分は、「$x$ を別の文字の式でおきかえる」向きにも使えます。半径 $r$ の円の面積を、積分で求めてみましょう。

\bkvpmidashi{円を細い長方形に分ける}
\noindent\begin{minipage}[c]{0.4\linewidth}
\centering
\begin{tikzpicture}[x=1mm, y=1mm]
  \fill[bkBlue!8] (0,0) circle (18);
  \draw[bkBlue, line width=1pt] (0,0) circle (18);
  \draw[bkGray, line width=0.5pt, -{Stealth[length=3pt]}] (-21,0) -- (22,0) node[anchor=west, font=\sffamily\scriptsize] {$x$};
  \draw[bkGray, line width=0.5pt, -{Stealth[length=3pt]}] (0,-21) -- (0,22) node[anchor=south, font=\sffamily\scriptsize] {$y$};
  \fill[bkTheme!35] (9.52,-14.74) rectangle (11.12,14.74);
  \draw[bkTheme!80!black, line width=0.5pt] (9.52,-14.74) rectangle (11.12,14.74);
  \draw[bkInk, line width=0.8pt] (0,0) -- (10.32,14.74);
  \draw[bkBlue!80!black, line width=1.8pt] (0,-0.1) -- (10.32,-0.1);
  \draw[bkOrange, line width=1.8pt] (12.3,0) -- (12.3,14.74);
  \draw[bkInk, line width=0.5pt] (0,6) arc (90:55:6);
  \node[font=\sffamily\scriptsize] at (72.5:8.6) {$\theta$};
  \node[font=\sffamily\scriptsize] at (6.9,6.0) {$r$};
  \fill[bkInk] (10.32,14.74) circle (0.8);
\end{tikzpicture}
\end{minipage}\hfill
\begin{minipage}[c]{0.56\linewidth}
円を、はば $\dd x$ のたての細長い長方形に分けます。位置 $x$ での長方形の上はしは円の上の点なので、中心からの距離は $r$。三平方の定理（4-4）から、上はしの高さ $y$ は
\[ x^2+y^2=r^2\quad\text{より}\quad y=\sqrt{r^2-x^2} \]
長方形は上下に同じだけあるので、高さは $2\sqrt{r^2-x^2}$。たとえば $r=5$、$x=3$ なら $y=\sqrt{25-9}=4$ で、高さ8です。
\end{minipage}\par\vspace{3pt}
この長方形を $x=-r$ から $r$ まで集めると、円の面積は
\[ S=\int_{-r}^{r}2\sqrt{r^2-x^2}\,\dd x \]
$\sqrt{\phantom{x}}$ の中がひき算で、このままでは積分できません。

\bkvpmidashi{図を見て、角 $\theta$ でおきかえる}
図の半径（長さ $r$）が、$y$ 軸となす角を $\theta$ とします。半径を斜めの辺とする直角三角形を見ると、
\begin{itemize}[leftmargin=1.2\zw, itemsep=1pt, topsep=2pt, label={\color{bkTheme}●}]
\item \textcolor{bkBlue!80!black}{横の長さ（青）} $x=r\sin\theta$
\item \textcolor{bkOrange!80!black}{たての長さ（オレンジ）} $\sqrt{r^2-x^2}=r\cos\theta$
\end{itemize}
（$\theta$ を $y$ 軸からはかったので、横が $\sin$、たてが $\cos$ になります。）だから、$x=r\sin\theta$ とおけば、$\sqrt{\phantom{x}}$ がはずれて $r\cos\theta$ になるはずです。式でも確かめます。$\sin^2\theta+\cos^2\theta=1$（4-8）から
\[ r^2-x^2=r^2-r^2\sin^2\theta=r^2(1-\sin^2\theta)=r^2\cos^2\theta,\qquad \sqrt{r^2\cos^2\theta}=r\cos\theta \]
（$-\frac{\pi}{2}\leqq\theta\leqq\frac{\pi}{2}$ では $\cos\theta\geqq0$ なので、$\sqrt{\phantom{x}}$ をはずしても $+$ のまま。）

\bkvpmidashi{$\dd x$ とはんいを書きかえる}
$x=r\sin\theta$ を $\theta$ で微分すると $\dfrac{\dd x}{\dd\theta}=r\cos\theta$（5-15）。だから $\dd x=r\cos\theta\,\dd\theta$ です。はんいは、図で半径を左はしから右はしまで動かして考えます。
\begin{bkcenter}
\begin{bktable}{colspec={Q[l,m,wd=30mm] *{3}{Q[c,m,wd=20mm]}}, rowsep=2pt}
 & 左はし & まん中 & 右はし \\
$x$ & $-r$ & 0 & $r$ \\
$\theta$ & $-\frac{\pi}{2}$ & 0 & $\frac{\pi}{2}$ \\
\end{bktable}
\end{bkcenter}
（$r\sin\bigl(-\frac{\pi}{2}\bigr)=-r$、$r\sin\frac{\pi}{2}=r$ で確かめられます。）

\bkvpmidashi{計算する}
\begin{align*}
S&=\int_{-\frac{\pi}{2}}^{\frac{\pi}{2}}\underbrace{2r\cos\theta}_{2\sqrt{r^2-x^2}}\cdot\underbrace{r\cos\theta\,\dd\theta}_{\dd x}=2r^2\int_{-\frac{\pi}{2}}^{\frac{\pi}{2}}\cos^2\theta\,\dd\theta=2r^2\int_{-\frac{\pi}{2}}^{\frac{\pi}{2}}\frac{1+\cos2\theta}{2}\,\dd\theta\\
&=2r^2\Bigl[\,\frac{\theta}{2}+\frac{\sin2\theta}{4}\,\Bigr]_{-\frac{\pi}{2}}^{\frac{\pi}{2}}=2r^2\Bigl\{\Bigl(\frac{\pi}{4}+\frac{\sin\pi}{4}\Bigr)-\Bigl(-\frac{\pi}{4}+\frac{\sin(-\pi)}{4}\Bigr)\Bigr\}\\
&=2r^2\Bigl(\frac{\pi}{4}+\frac{\pi}{4}\Bigr)=2r^2\times\frac{\pi}{2}=\pi r^2
\end{align*}
（$\cos^2\theta=\dfrac{1+\cos2\theta}{2}$ は5-15。$\sin\pi=\sin(-\pi)=0$。）小学校から使ってきた「円の面積＝半径×半径×円周率」が、積分で確かめられました。

\bksubsection{7}{積の微分を、逆に読む：部分積分}

\label{bk:5-16:bubun}こんどは、積の微分を逆に読みます。まず、数の式で見てみましょう。$x\,e^x$ を積の微分（5-6）で微分すると、
\[ (x\,e^x)'=\underbrace{(x)'e^x}_{1\cdot e^x}+\underbrace{x\,(e^x)'}_{x\,e^x}=e^x+x\,e^x \]
これを逆に読むと、「$e^x+x\,e^x$ を積分すると $x\,e^x$」。2つのたし算を1つずつ積分すると、
\[ \int e^x\,\dd x+\int x\,e^x\,\dd x=x\,e^x \]
（$+C$ は最後にまとめてつけます。）$\int e^x\,\dd x$ を右へ移すと（$\int e^x\,\dd x=e^x$ は5-14）、
\[ \int x\,e^x\,\dd x=x\,e^x-\int e^x\,\dd x=x\,e^x-e^x+C \]
表になかった $\int x\,e^x\,\dd x$ が求められました。微分して確かめると、
\[ (xe^x-e^x)'=(e^x+xe^x)-e^x=xe^x \]
で、合っています。

\bkvpmidashi{文字で書くと}
同じことを、$f$ と $g$ で書きます。積の微分 $(fg)'=f'g+fg'$ の両辺を積分すると、左は「微分したものを積分する」のでもとの $fg$ にもどります（5-11）。
\[ f\,g=\int f'g\,\dd x+\int f\,g'\,\dd x \]
$\int f'g\,\dd x$ を左へ移して、左右を入れかえると、
\[ \int f\,g'\,\dd x=f\,g-\int f'\,g\,\dd x \]
これを\textbf{部分積分}といいます。左の積分が求めにくくても、右の $\int f'g\,\dd x$ が簡単になれば、答えが出ます。上の例は、$f=x$、$g'=e^x$（$g=e^x$）としたものです。\bkidx{ぶぶんせきぶん}{部分積分}

\bkvpmidashi{使い方の手順}
\begin{enumerate}[label={\bkcircled[bkTheme]{\arabic*}}, leftmargin=2.4\zw, itemsep=1pt, topsep=2pt]
\item かけ算の2つのうち、どちらを $f$（微分するほう）、どちらを $g'$（積分するほう）にするか決める（⑧）。
\item $f'$（$f$ を微分したもの）と、$g$（$g'$ を積分したもの。$+C$ はつけない）を求める。
\item $fg-\int f'g\,\dd x$ に入れて、残った積分を計算する。最後に $+C$ をつける。
\end{enumerate}

\begin{bkfigure}
\begin{tikzpicture}[x=1mm, y=1mm]
  \fill[bkBlue!14] (0,0) -- (0,24) -- plot[domain=0:1, samples=30] ({50*\x}, {24+12*\x*\x}) -- (50,0) -- cycle;
  \fill[bkOrange!20] (0,24) -- plot[domain=0:1, samples=30] ({50*\x}, {24+12*\x*\x}) -- (0,36) -- cycle;
  \draw[bkInk, line width=1.1pt] plot[domain=0:1, samples=30] ({50*\x}, {24+12*\x*\x});
  \draw[bkInk, line width=0.8pt] (0,0) rectangle (50,36);
  \node[font=\small, text=bkBlue!70!black] at (25,12) {$\int f\,\dd g$};
  \node[font=\small, text=bkOrange!80!black] at (12,31) {$\int g\,\dd f$};
  \node[font=\sffamily\scriptsize, anchor=north] at (25,-1) {横：$g$};
  \node[font=\sffamily\scriptsize, anchor=east] at (-1,18) {たて：$f$};
  \node[font=\sffamily\small, anchor=west, align=left] at (56,18) {全体の長方形＝$f\,g$\\青＋オレンジ＝$f\,g$\\だから\enspace 青＝$f\,g-{}$オレンジ};
\end{tikzpicture}
\bkfigcaption{図で見る部分積分。横に $g$、たてに $f$ をとって、$f$ と $g$ が変わっていくようすを曲線でかく。青い部分は、はば $\dd g$・高さ $f$ のたての細い長方形を集めたもの。オレンジの部分は、はば $\dd f$・長さ $g$ の横の細い長方形を集めたもの。$g'\,\dd x=\dd g$、$f'\,\dd x=\dd f$ なので、部分積分は「青＝全体${}-{}$オレンジ」（はじめ $g=0$ の場合）。}
\end{bkfigure}

\bksubsection{8}{どちらを $f$ にするか}

\label{bk:5-16:erabu}部分積分では、\textbf{微分すると簡単になるもの}を $f$、\textbf{積分しやすいもの}を $g'$ にします。$\displaystyle\int x\,e^x\,\dd x$ で、逆に $f=e^x$、$g'=x$（$g=\frac{x^2}{2}$）とすると、
\[ \int x\,e^x\,\dd x=\frac{x^2}{2}e^x-\int\frac{x^2}{2}e^x\,\dd x \]
となり、残った積分に $x^2$ が出てきて、はじめよりむずかしくなります。$x$ は微分すると1になって消えるので、$f$ にするのがよいのです。よく出る形の選び方は、次の表のとおりです。
\begin{bkcenter}
\begin{tblr}{colspec={Q[l,m]Q[l,m]Q[l,m]}, row{1}={bg=bkTheme!12, font=\sffamily\small}, hline{1,2,Z}={0.6pt,bkGray}, hline{3-5}={0.3pt,bkGray}, colsep=6pt}
積分の形 & $f$（微分する） & $g'$（積分する） \\
$x\,e^{ax}$ & $x$ → 1 & $e^{ax}$ → $\frac{1}{a}e^{ax}$ \\
$x\sin x$、$x\cos x$ & $x$ → 1 & $\sin x$ → $-\cos x$、$\cos x$ → $\sin x$ \\
$x^2e^x$ & $x^2$ → $2x$（2回使う） & $e^x$ → $e^x$ \\
$\ln x$ & $\ln x$ → $\frac{1}{x}$ & 1 → $x$ \\
\end{tblr}
\end{bkcenter}

\bkvpmidashi{$\ln x$ の積分}
$\ln x$ は1つの関数だけで、かけ算に見えません。でも、$\ln x=\ln x\times1$ と見ると、部分積分が使えます。$\ln x$ は積分のしかたが分かりませんが、微分すると $\frac{1}{x}$ と簡単になるので、$f=\ln x$、$g'=1$（$g=x$）とします。
\[ \int\ln x\,\dd x=\underbrace{\ln x}_{f}\cdot\underbrace{x}_{g}-\int\underbrace{\frac{1}{x}}_{f'}\cdot\underbrace{x}_{g}\,\dd x=x\ln x-\int1\,\dd x=x\ln x-x+C \]
確かめ：$(x\ln x-x)'=\Bigl(1\cdot\ln x+x\cdot\dfrac{1}{x}\Bigr)-1=\ln x+1-1=\ln x$。合っています。

\bksubsection{9}{定積分での部分積分}

\label{bk:5-16:teisekibun}定積分では、$fg$ の部分にも、はんいをつけます。
\[ \int_a^bf\,g'\,\dd x=\Bigl[\,f\,g\,\Bigr]_a^b-\int_a^bf'\,g\,\dd x \]
$\Bigl[\,f\,g\,\Bigr]_a^b$ は、$f\,g$ の $b$ での値から $a$ での値をひいたものです（5-11）。

\bksubsection{10}{はてな？の答え}

\label{bk:5-16:kotae}\bkhatenaref の電気の量 $\displaystyle\int_0^32t\,e^{-t}\,\dd t$ を求めましょう。
\begin{enumerate}[label={\bkcircled[bkTheme]{\arabic*}}, leftmargin=2.4\zw, itemsep=2pt, topsep=2pt]
\item $2t$ は微分すると2と簡単になるので $f=2t$。$e^{-t}$ を $g'$ にします。
\item $f'=2$。$g$ は $e^{-t}$ を積分したもので、5-14の $\int e^{at}\,\dd t=\frac{1}{a}e^{at}$ で $a=-1$ とすると、$g=-e^{-t}$。
\item 部分積分に入れます。$f'g=2\times(-e^{-t})=-2e^{-t}$ なので、ひき算の $-$ と合わせて $+$ になります。
\begin{align*}
\int_0^32t\,e^{-t}\,\dd t&=\Bigl[\,2t\times(-e^{-t})\,\Bigr]_0^3-\int_0^3(-2e^{-t})\,\dd t\\
&=\Bigl[\,-2t\,e^{-t}\,\Bigr]_0^3+\int_0^32e^{-t}\,\dd t
\end{align*}
\item 前の部分：$t=3$ で $-6e^{-3}$、$t=0$ で $0$。ひいて $-6e^{-3}$。
\item 後ろの部分：$2e^{-t}$ の原始関数は $-2e^{-t}$ なので、
\[ \Bigl[\,-2e^{-t}\,\Bigr]_0^3=-2e^{-3}-(-2e^0)=-2e^{-3}+2 \]
\item 合わせて $-6e^{-3}-2e^{-3}+2=2-8e^{-3}$。
\end{enumerate}
$e^{-3}=0.0498$ なので、$2-8\times0.0498=2-0.398=1.602$。答えは約 \qty{1.60}{C} です。表計算の1.6017と合いました。

\bksubsection{11}{積分では、どの道具を使う？}

\label{bk:5-16:erabikata}積分の式を見たら、次の順に考えます。微分の章（5-6）の「どの道具を使う？」と同じように、上から順に問いかけます。
\begin{bkfigure}
\begin{tikzpicture}[x=1mm, y=1mm, q/.style={bk node, minimum width=50mm, minimum height=13mm, font=\small},
  r/.style={bk node, draw=bkOrange, fill=bkOrange!8, minimum width=46mm, minimum height=13mm, font=\small}]
  \node[q] (q1) at (25,0) {積分の表（5-11・5-14・5-15）に\\そのままある？};
  \node[r] (a1) at (97,0) {表を逆に読んで積分};
  \node[q] (q2) at (25,-22) {「内側の式」と\\「内側の微分」が入っている？};
  \node[r] (a2) at (97,-22) {置換積分（②〜⑥）};
  \node[q] (q3) at (25,-44) {$x$ や $\ln x$ と、$e^x$・$\sin$・$\cos$ の\\かけ算になっている？};
  \node[r] (a3) at (97,-44) {部分積分（⑦〜⑩）};
  \node[r, minimum width=50mm] (a4) at (25,-66) {式を変形してから考えなおす\\（2倍角の式など、5-15）};
  \draw[bk link] (q1) -- node[bk link label, above] {はい} (a1);
  \draw[bk link] (q1) -- node[bk link label, right] {いいえ} (q2);
  \draw[bk link] (q2) -- node[bk link label, above] {はい} (a2);
  \draw[bk link] (q2) -- node[bk link label, right] {いいえ} (q3);
  \draw[bk link] (q3) -- node[bk link label, above] {はい} (a3);
  \draw[bk link] (q3) -- node[bk link label, right] {いいえ} (a4);
\end{tikzpicture}
\bkfigcaption{積分の道具の選び方。どの道具を使っても、さいごは答えを微分して、もとの式にもどるか確かめる。}
\end{bkfigure}
たとえば $\displaystyle\int\sin^2t\,\dd t$ は、表にもなく、内側の微分もなく、かけ算の形でもないので、2倍角の式 $\sin^2t=\dfrac{1-\cos2t}{2}$ で形を変えてから積分しました（5-15）。

\begin{motto}[4]
\textbf{2回使うと、もとにもどる}\par
$I=\displaystyle\int e^x\sin x\,\dd x$ に部分積分を2回使います。1回めは $f=\sin x$、$g'=e^x$。2回めは $f=\cos x$、$g'=e^x$ です。
\begin{align*}
I&=e^x\sin x-\int e^x\cos x\,\dd x\\
&=e^x\sin x-\Bigl(e^x\cos x+\int e^x\sin x\,\dd x\Bigr)=e^x(\sin x-\cos x)-I
\end{align*}
（2行めのかっこの中は、$\int e^x\cos x\,\dd x=e^x\cos x-\int(-\sin x)e^x\,\dd x$ です。）右にも $I$ が出てきました。$I$ を左に移すと $2I=e^x(\sin x-\cos x)$、つまり $I=\dfrac{1}{2}e^x(\sin x-\cos x)+C$。ゆれながら止まる動き（減衰振動、6-7・6-12）の計算で使います。
\end{motto}

\begin{senmonsho}[大学の教科書ふうの書き方]
\begin{bkhonbun}
部分積分の公式 $\int u\,\dd v=uv-\int v\,\dd u$ を用いると、次式を得る。
\[ \int_0^\infty e^{-st}\frac{\dd f}{\dd t}\,\dd t=\Bigl[\,e^{-st}f(t)\,\Bigr]_0^\infty+s\int_0^\infty e^{-st}f(t)\,\dd t \]
\end{bkhonbun}
\bkyaku{$u$ と $v$ は、この章の $f$ と $g$ のこと。$\dd v=g'\,\dd x$。$\int_0^\infty$ は「0からどこまでも」積分すること。$e^{-st}f'(t)$ の積分が、$f(t)$ の積分の $s$ 倍に変わる。これが9-5のラプラス変換で「微分がかけ算に変わる」しくみ。}
\end{senmonsho}

\begin{taisetsu}
\def\tsl#1{\par\addvspace{3pt}\noindent\hangindent=5.2\zw\hangafter=1\makebox[5.2\zw][l]{\sffamily\bfseries\small\color{bkTheme}#1}}
\tsl{置換積分}中身 $u$ とその微分を見つけ、$\dd u=u'\,\dd x$ で書きかえる。定積分は、はんいも $u$ に。
\tsl{部分積分}$\displaystyle\int f\,g'\,\dd x=f\,g-\int f'\,g\,\dd x$\par
\noindent\hspace*{5.2\zw}$f$：微分で簡単になるもの、$g'$：積分しやすいもの
\tsl{確かめ}答えを微分して、積分の中身にもどるか確かめる。
\end{taisetsu}

\begin{kotoba}
\small
\bkword{置換積分}{積分の中の式を $u$ におきかえて、簡単な形にする方法。合成関数の微分の逆。}
\bkword{部分積分}{積の微分の逆。$\int fg'\,\dd x=fg-\int f'g\,\dd x$。}
\end{kotoba}

%% ---------------- 例題 ----------------
\bksection{reidai}{例題でたしかめよう}

\begin{reidai}[置換積分]
$\displaystyle\int3x^2(x^3+2)^4\,\dd x$ を求めましょう。\par
\kangaekata かっこの中 $x^3+2$ を微分すると $3x^2$ で、それが前にある。だから $u=x^3+2$ とおく（③）。\par
\shiki $u=x^3+2$ とおくと $\dfrac{\dd u}{\dd x}=3x^2$ より $\dd u=3x^2\,\dd x$\par
\hspace*{2.6\zw}$\displaystyle\int\underbrace{(x^3+2)^4}_{u^4}\,\underbrace{3x^2\,\dd x}_{\dd u}=\int u^4\,\dd u=\frac{u^5}{5}+C$\par
\hspace*{2.6\zw}$u$ を $x^3+2$ にもどす。確かめ：\par
\hspace*{4.6\zw}$\Bigl\{\dfrac{(x^3+2)^5}{5}\Bigr\}'=\dfrac{5(x^3+2)^4\times3x^2}{5}=3x^2(x^3+2)^4$\par
\kotae $\dfrac{(x^3+2)^5}{5}+C$
\end{reidai}
置換積分では、さいごに $u$ を $x$ の式にもどすのをわすれないようにしましょう。問題は $x$ の式なので、答えも $x$ の式で書きます。もどしたら、微分して積分の中身にもどるかを確かめると安心です。

\begin{jibun}
$\displaystyle\int2x(x^2-1)^5\,\dd x$ を求めましょう。
\bkjans{$u=x^2-1$、$\dd u=2x\,\dd x$。$\dfrac{(x^2-1)^6}{6}+C$}
\end{jibun}

\begin{reidai}[係数を調整する置換積分]
（1）$\displaystyle\int t\,e^{-t^2}\,\dd t$\quad（2）$\displaystyle\int\cos x\,\sin^3x\,\dd x$\quad（3）$\displaystyle\int\sin x\,\cos^3x\,\dd x$ を求めましょう。\par
\kangaekata（1）$u=-t^2$ とおき、数のずれを④のように調整。（2）$\cos x$ は $\sin x$ の微分なので $u=\sin x$。（3）$\cos x$ の微分は $-\sin x$ なので $u=\cos x$ とし、$-$ を調整。\par
\shiki（1）$\dd u=-2t\,\dd t$ より $t\,\dd t=-\dfrac{1}{2}\dd u$。$\displaystyle\int e^u\Bigl(-\frac{1}{2}\Bigr)\dd u=-\frac{1}{2}e^u+C$\par
\hspace*{2.6\zw}（2）$\dd u=\cos x\,\dd x$。$\displaystyle\int u^3\,\dd u=\frac{u^4}{4}+C$\par
\hspace*{2.6\zw}（3）$\dd u=-\sin x\,\dd x$ より $\sin x\,\dd x=-\dd u$。\par
\hspace*{2.6\zw}\phantom{（3）}$\displaystyle\int u^3(-1)\,\dd u=-\frac{u^4}{4}+C$\par
\kotae （1）$-\dfrac{1}{2}e^{-t^2}+C$\quad（2）$\dfrac{\sin^4x}{4}+C$\quad（3）$-\dfrac{\cos^4x}{4}+C$
\end{reidai}
（3）を微分して確かめると、
\[ \Bigl(-\frac{\cos^4x}{4}\Bigr)'=-\frac{4\cos^3x\times(-\sin x)}{4}=\sin x\cos^3x \]
$-$ が2つで $+$ になって、もとにもどります。

\begin{jibun}
（1）$\displaystyle\int x\cos(x^2)\,\dd x$\quad（2）$\displaystyle\int\sin x\,\cos^2x\,\dd x$ を求めましょう。
\bkjans{（1）\nobreak $u=x^2$、$x\,\dd x=\frac{1}{2}\dd u$。$\dfrac{1}{2}\sin(x^2)+C$\quad（2）\nobreak $u=\cos x$、$\sin x\,\dd x=-\dd u$。$-\dfrac{\cos^3x}{3}+C$}
\end{jibun}

\begin{reidai}[分母の微分が分子にある形]
（1）$\displaystyle\int\frac{e^x}{e^x+1}\,\dd x$\quad（2）$\displaystyle\int\frac{x}{x^2+4}\,\dd x$ を求めましょう。\par
\kangaekata 分子が分母の微分（またはその数倍）になっている。④の「よく出る形」で $u=$（分母）。\par
\shiki（1）$u=e^x+1$（いつも正）とおくと $\dd u=e^x\,\dd x$。$\displaystyle\int\frac{1}{u}\,\dd u=\ln u+C$\par
\hspace*{2.6\zw}（2）$u=x^2+4$（いつも正）とおくと $\dd u=2x\,\dd x$ より $x\,\dd x=\dfrac{1}{2}\dd u$\par
\hspace*{4.6\zw}$\displaystyle\int\frac{1}{u}\cdot\frac{1}{2}\,\dd u=\frac{1}{2}\ln u+C$\par
\kotae （1）$\ln(e^x+1)+C$\quad（2）$\dfrac{1}{2}\ln(x^2+4)+C$
\end{reidai}
「いつも正」をたしかめるのは、$\ln u$ は $u>0$ のときだけ考えるからです（5-14）。$e^x$ も $x^2$ も0以上なので、$e^x+1$ や $x^2+4$ はいつも正です。（2）のように、分子が分母の微分のちょうど何倍か（ここでは $\frac{1}{2}$ 倍）になっているときは、その数を前に出してから積分します。

\begin{jibun}
$\displaystyle\int\frac{4x}{x^2+3}\,\dd x$ を求めましょう。
\bkjans{$u=x^2+3$、$\dd u=2x\,\dd x$ より $4x\,\dd x=2\,\dd u$。$2\ln(x^2+3)+C$}
\end{jibun}

\begin{reidai}[定積分の置換]
（1）$\displaystyle\int_0^1x(x^2+1)^2\,\dd x$\quad（2）$\displaystyle\int_0^{\frac{\pi}{2}}\sin^2t\,\cos t\,\dd t$ を求めましょう。\par
\kangaekata $u$ におきかえ、はんいも $u$ の値に書きかえる（⑤）。\par
\shiki（1）$u=x^2+1$、$x\,\dd x=\dfrac{1}{2}\dd u$。$x:0\to1$ のとき $u:1\to2$\par\vspace{2pt}
\hspace*{4.6\zw}$\displaystyle\frac{1}{2}\int_1^2u^2\,\dd u=\frac{1}{2}\Bigl[\,\frac{u^3}{3}\,\Bigr]_1^2=\frac{1}{2}\times\Bigl(\frac{8}{3}-\frac{1}{3}\Bigr)=\frac{1}{2}\times\frac{7}{3}=\frac{7}{6}$\par
\hspace*{2.6\zw}（2）$u=\sin t$、$\dd u=\cos t\,\dd t$。$t:0\to\frac{\pi}{2}$ のとき $u:\sin0=0\to\sin\frac{\pi}{2}=1$\par\vspace{2pt}
\hspace*{4.6\zw}$\displaystyle\int_0^1u^2\,\dd u=\Bigl[\,\frac{u^3}{3}\,\Bigr]_0^1=\frac{1}{3}-0=\frac{1}{3}$\par
\kotae （1）$\dfrac{7}{6}$\quad（2）$\dfrac{1}{3}$
\end{reidai}
（2）では、$t$ が0から $\frac{\pi}{2}$ まで動くと、$u=\sin t$ は0から1まで動きます。はんいの書きかえは、下端と上端のそれぞれに、おいた式（ここでは $\sin t$）を計算するだけです。（1）は、$x$ の式にもどしてから計算しても同じです。
\[ \Bigl[\,\frac{(x^2+1)^3}{6}\,\Bigr]_0^1=\frac{8}{6}-\frac{1}{6}=\frac{7}{6} \]

\begin{jibun}
$\displaystyle\int_0^12x(x^2+1)^3\,\dd x$ と $\displaystyle\int_0^{\frac{\pi}{2}}\sin t\,\cos t\,\dd t$ を求めましょう。
\bkjans{$u=x^2+1$（$1\to2$）：$\dfrac{15}{4}$。$u=\sin t$（$0\to1$）：$\Bigl[\,\dfrac{u^2}{2}\,\Bigr]_0^1=\dfrac{1}{2}$}
\end{jibun}

\begin{reidai}[$x=r\sin\theta$ とおく置換]
$\displaystyle\int_0^2\sqrt{4-x^2}\,\dd x$ を求めましょう。また、それは何の面積ですか。\par
\kangaekata ⑥の円の面積と同じく、$r=2$ として $x=2\sin\theta$ とおき、$\sqrt{\phantom{x}}$ をはずす。\par
\shiki $x=2\sin\theta$ とおくと\par
\hspace*{4.6\zw}$\sqrt{4-x^2}=\sqrt{4-4\sin^2\theta}=\sqrt{4\cos^2\theta}=2\cos\theta$\par
\hspace*{2.6\zw}$\dd x=2\cos\theta\,\dd\theta$。$x:0\to2$ のとき $\theta:0\to\frac{\pi}{2}$（$2\sin\frac{\pi}{2}=2$）\par\vspace{3pt}
\hspace*{2.6\zw}$\displaystyle\int_0^{\frac{\pi}{2}}2\cos\theta\cdot2\cos\theta\,\dd\theta=4\int_0^{\frac{\pi}{2}}\frac{1+\cos2\theta}{2}\,\dd\theta$\par\vspace{3pt}
\hspace*{4.6\zw}$\displaystyle=4\Bigl[\,\frac{\theta}{2}+\frac{\sin2\theta}{4}\,\Bigr]_0^{\frac{\pi}{2}}=4\times\Bigl(\frac{\pi}{4}+0\Bigr)-0=\pi$\par
\kotae $\pi$（半径2の円の $\frac{1}{4}$ の面積。$\pi\times2^2\div4=\pi$ と合う）
\end{reidai}
\noindent\begin{minipage}[c]{0.3\linewidth}
\centering
\begin{tikzpicture}[x=1mm, y=1mm]
  \fill[bkBlue!15] (0,0) -- (20,0) arc (0:90:20) -- cycle;
  \draw[bkBlue, line width=1pt] (20,0) arc (0:90:20);
  \draw[bkGray, line width=0.5pt, -{Stealth[length=3pt]}] (-2,0) -- (25,0) node[anchor=west, font=\sffamily\scriptsize] {$x$};
  \draw[bkGray, line width=0.5pt, -{Stealth[length=3pt]}] (0,-2) -- (0,25) node[anchor=south, font=\sffamily\scriptsize] {$y$};
  \node[font=\sffamily\scriptsize, anchor=north] at (20,-1) {2};
  \node[font=\sffamily\scriptsize, anchor=east] at (-1,20) {2};
\end{tikzpicture}
\end{minipage}\hfill
\begin{minipage}[c]{0.66\linewidth}
$y=\sqrt{4-x^2}$ は、両辺を2乗すると $x^2+y^2=4$ で、半径2の円の上半分です。0から2までなので、左の図の青い部分（円の $\frac{1}{4}$）の面積を求めたことになります。
\end{minipage}

\begin{jibun}
$\displaystyle\int_0^3\sqrt{9-x^2}\,\dd x$ を求めましょう。
\bkjans{$x=3\sin\theta$、$\dd x=3\cos\theta\,\dd\theta$。$\displaystyle\int_0^{\frac{\pi}{2}}9\cos^2\theta\,\dd\theta=\dfrac{9\pi}{4}$（半径3の円の $\frac{1}{4}$）}
\end{jibun}

\begin{reidai}[部分積分]
（1）$\displaystyle\int x\cos x\,\dd x$\quad（2）$\displaystyle\int x\,e^{2x}\,\dd x$ を求めましょう。\par
\kangaekata $x$ とほかの関数の積なので、⑦の部分積分。$f$ は微分すると1になる $x$（⑧）。\par
\shiki（1）$f=x$、$g'=\cos x$ とする。$f'=1$、$g=\sin x$\par
\hspace*{4.6\zw}$\displaystyle\int x\cos x\,\dd x=x\sin x-\int1\cdot\sin x\,\dd x$\par
\hspace*{4.6\zw}$\displaystyle\phantom{\int x\cos x\,\dd x}=x\sin x-(-\cos x)+C$\par
\hspace*{2.6\zw}（2）$f=x$、$g'=e^{2x}$ とする。$f'=1$、$g=\dfrac{1}{2}e^{2x}$（5-14）\par
\hspace*{4.6\zw}$\displaystyle\int x\,e^{2x}\,\dd x=\frac{1}{2}x\,e^{2x}-\int\frac{1}{2}e^{2x}\,\dd x=\frac{1}{2}x\,e^{2x}-\frac{1}{4}e^{2x}+C$\par
\kotae （1）$x\sin x+\cos x+C$\quad（2）$\dfrac{1}{2}xe^{2x}-\dfrac{1}{4}e^{2x}+C$
\end{reidai}
（1）を微分して確かめます。$x\sin x$ は積の微分（5-6）で微分します。
\[ (x\sin x+\cos x)'=(1\cdot\sin x+x\cos x)-\sin x=x\cos x \]
もとの式にもどりました。

\begin{jibun}
（1）$\displaystyle\int x\,e^{-x}\,\dd x$\quad（2）$\displaystyle\int x\sin x\,\dd x$ を求めましょう。
\bkjans{（1）\nobreak $g=-e^{-x}$。$-xe^{-x}-e^{-x}+C$\quad（2）\nobreak $g=-\cos x$。$-x\cos x+\sin x+C$}
\end{jibun}

\begin{reidai}[$\ln x$ の積分]
$\displaystyle\int_1^e\ln x\,\dd x$ を求めましょう。\par
\kangaekata ⑧の $\ln x$ の積分 $x\ln x-x$ を使い、上端の値から下端の値をひく。\par
\shiki $\ln e=1$（$e^1=e$）、$\ln1=0$（$e^0=1$）より\par
\hspace*{2.6\zw}$x=e$ で $e\cdot1-e=0$、$x=1$ で $1\cdot0-1=-1$\par
\hspace*{2.6\zw}$\Bigl[\,x\ln x-x\,\Bigr]_1^e=0-(-1)=1$\par
\kotae 1
\end{reidai}

\begin{jibun}
$\displaystyle\int_1^2\ln x\,\dd x$ を求めましょう（$\ln2=0.693$）。
\bkjans{$\Bigl[\,x\ln x-x\,\Bigr]_1^2=(2\ln2-2)-(0-1)=2\ln2-1\approx0.386$}
\end{jibun}

\begin{reidai}[定積分での部分積分：流れた電気の量]
電流 $i=2t\,e^{-t}$（A）が流れました。はじめの3秒間に流れた電気の量を求めましょう（$e^{-3}=0.0498$）。\par
\kangaekata 電気の量は電流の積分（5-12）。積なので⑨の部分積分（⑩と同じ手順）。\par
\shiki $f=2t$、$g'=e^{-t}$。$f'=2$、$g=-e^{-t}$\par
\hspace*{2.6\zw}$\begin{aligned}[t]\int_0^32t\,e^{-t}\,\dd t&=\Bigl[\,-2te^{-t}\,\Bigr]_0^3+\int_0^32e^{-t}\,\dd t=-6e^{-3}+\Bigl[\,-2e^{-t}\,\Bigr]_0^3\\
&=-6e^{-3}-2e^{-3}+2=2-8e^{-3}=2-8\times0.0498\end{aligned}$\par
\kotae $2-8e^{-3}\approx\qty{1.60}{C}$
\end{reidai}

\begin{jibun}
$\displaystyle\int_0^{\pi}t\sin t\,\dd t$ を求めましょう。
\bkjans{$f=t$、$g=-\cos t$。$\Bigl[\,-t\cos t\,\Bigr]_0^{\pi}+\displaystyle\int_0^{\pi}\cos t\,\dd t=\pi+0=\pi$}
\end{jibun}

\begin{reidai}[部分積分を2回使う]
$\displaystyle\int x^2e^x\,\dd x$ を求めましょう。\par
\kangaekata ⑧の表の $x^2e^x$ の行。1回めで $x^2$ が $2x$ になり、2回めで定数になる。\par
\shiki 1回め：$f=x^2$、$g'=e^x$。$f'=2x$、$g=e^x$\par
\hspace*{4.6\zw}$\displaystyle\int x^2e^x\,\dd x=x^2e^x-\int2x\,e^x\,\dd x$\par
\hspace*{2.6\zw}2回め：⑦で求めた $\displaystyle\int x\,e^x\,\dd x=xe^x-e^x$ を使う\par
\hspace*{4.6\zw}$=x^2e^x-2(xe^x-e^x)+C=x^2e^x-2xe^x+2e^x+C$\par
\kotae $e^x(x^2-2x+2)+C$
\end{reidai}

\begin{jibun}
$\displaystyle\int_0^1x^2e^x\,\dd x$ を求めましょう（$e=2.718$）。
\bkjans{$\Bigl[\,e^x(x^2-2x+2)\,\Bigr]_0^1=e\times1-1\times2=e-2\approx0.718$}
\end{jibun}

%% ---------------- まちがえやすいところ・物理工学 ----------------
\bksection{tsumazuki}{まちがえやすいところをチェック}

\begin{tsumazuki}{定積分の置換で、はんいを書きかえない}
\bkng{$\displaystyle\int_0^12x(x^2+1)^3\,\dd x=\int_0^1u^3\,\dd u$\quad（はんいが $x$ のまま）}
\bkok{$\displaystyle\int_0^12x(x^2+1)^3\,\dd x=\int_1^2u^3\,\dd u$}
はんいは $x$ の値のままではなく、$u=x^2+1$ の値（1と2）に書きかえます。×のまま計算すると $\frac{1}{4}$ になり、正しい $\frac{15}{4}$ とずれます。
\end{tsumazuki}

\begin{tsumazuki}{$\dd x$ を $\dd u$ に書きかえわすれる}
$\displaystyle\int x\,e^{x^2}\,\dd x$ で、$e^{x^2}$ を $e^u$ にしただけで $\int e^u\,\dd x$ としてはいけません。$x\,\dd x=\frac{1}{2}\dd u$ まで書きかえて、積分の中を $u$ だけにします。
\end{tsumazuki}

\begin{tsumazuki}{部分積分の $-$ をわすれる}
\bkng{$\displaystyle\int x\cos x\,\dd x=x\sin x+\int\sin x\,\dd x$}
\bkok{$\displaystyle\int x\cos x\,\dd x=x\sin x-\int\sin x\,\dd x$}
部分積分は「$fg$ ひく $\int f'g$」です。答えを微分して、もとの式にもどるか確かめましょう。計算が進むほど式が長くなったら、$f$ と $g'$ の選び方も見直します（⑧）。
\end{tsumazuki}

\begin{butsuri}
\begin{itemize}[leftmargin=1.2\zw, itemsep=1pt, topsep=0pt, label={\color{bkPurple}●}]
\item \begin{minipage}[t]{0.53\linewidth}\raggedright
\par\addvspace{3pt}\textbf{流れた電気の量}\par\noindent スイッチを入れたあとの電流は、$t\,e^{-at}$ の形のことが多くあります。$i=2te^{-t}$ を、0から $t$ 秒まで部分積分すると、$t$ 秒までに流れた電気の量は\par\vspace{3pt}
{\centering $Q=2-2(t+1)e^{-t}$\par}\vspace{3pt}
となり、2Cに近づきます。
\end{minipage}\hfill
\begin{minipage}[t]{0.45\linewidth}
\begin{tikzpicture}[baseline={([yshift=-3.5mm]current bounding box.north)}]
\begin{axis}[bkplot, scale only axis, width=36mm, height=27mm, xmin=0, xmax=6, ymin=0, ymax=2.4, xtick={1,2,...,5}, ytick={1,2}, xlabel={$t$（秒）}]
  \addplot[bkcurve, domain=0:5.6, samples=80] {2*x*exp(-x)};
  \addplot[bkcurve2, domain=0:5.4, samples=80] {2-2*(x+1)*exp(-x)};
  \node[font=\sffamily\footnotesize, text=bkBlue, anchor=south west, fill=bkPurpleBg, inner sep=1pt] at (axis cs:3.0,0.45) {電流 $i$（A）};
  \node[font=\sffamily\footnotesize, text=bkOrange!80!black, anchor=north west, align=left, inner sep=1pt] at (axis cs:6.1,2.0) {流れた量\\[-2pt]$Q$（C）};
\end{axis}
\end{tikzpicture}
\end{minipage}\par\vspace{9pt}
\item \textbf{ラプラス変換（9-5）}\par\noindent 制御工学の中心の道具。$e^{-st}$ をかけて積分すると、微分が $s$ のかけ算に変わります。そのしくみが、部分積分です。
\item \textbf{フーリエ級数（9-1）}\par\noindent $\int t\sin nt\,\dd t$ のような積分を部分積分で計算して、音や振動を $\sin$ と $\cos$ に分けます。
\end{itemize}
\end{butsuri}

%% ---------------- 章末問題 ----------------
\bkendproblems
\bkproblemgroup{基本}{3}

\bkproblem 次の不定積分を求めましょう。\quad（1）$\displaystyle\int4x(x^2+5)^3\,\dd x$\quad（2）$\displaystyle\int x\,e^{3x^2}\,\dd x$\quad（3）$\displaystyle\int\cos x\,\sin^2x\,\dd x$
\bkpans{（1）\nobreak $\dfrac{(x^2+5)^4}{2}+C$\quad（2）\nobreak $\dfrac{1}{6}e^{3x^2}+C$\quad（3）\nobreak $\dfrac{\sin^3x}{3}+C$}

\bkproblem $\displaystyle\int\frac{2x+1}{x^2+x+1}\,\dd x$ を求めましょう（$x^2+x+1$ はいつも正）。
\bkpans{$u=x^2+x+1$。$\ln(x^2+x+1)+C$}

\bkproblem 次の定積分を求めましょう。\quad（1）$\displaystyle\int_0^2x(x^2+1)^2\,\dd x$\quad（2）$\displaystyle\int_0^1t\,e^{t^2}\,\dd t$
\bkpans{（1）\nobreak $\dfrac{1}{2}\Bigl[\,\dfrac{u^3}{3}\,\Bigr]_1^5=\dfrac{62}{3}$\quad（2）\nobreak $\dfrac{1}{2}\Bigl[\,e^u\,\Bigr]_0^1=\dfrac{e-1}{2}\approx0.859$}

\bkproblem 次の不定積分を求めましょう。\quad（1）$\displaystyle\int x\,e^{2x}\,\dd x$\quad（2）$\displaystyle\int3x\cos x\,\dd x$
\bkpans{（1）\nobreak $\dfrac{1}{2}xe^{2x}-\dfrac{1}{4}e^{2x}+C$\quad（2）\nobreak $3x\sin x+3\cos x+C$}

\bkproblem $\displaystyle\int_0^{\frac{\pi}{2}}x\cos x\,\dd x$ と $\displaystyle\int_1^e2\ln x\,\dd x$ を求めましょう。
\bkpans{$\Bigl[\,x\sin x+\cos x\,\Bigr]_0^{\frac{\pi}{2}}=\dfrac{\pi}{2}-1\approx0.571$。$2\times1=2$}

\bkproblem 次を求めましょう。\quad（1）$\displaystyle\int3x^2(x^3+1)^2\,\dd x$\quad（2）$\displaystyle\int x\sin2x\,\dd x$
\bkpans{（1）\nobreak $u=x^3+1$。$\dfrac{(x^3+1)^3}{3}+C$\quad（2）\nobreak $f=x$、$g=-\dfrac{1}{2}\cos2x$。$-\dfrac{1}{2}x\cos2x+\dfrac{1}{4}\sin2x+C$}

\bkproblem $\displaystyle\int_0^1x\,e^x\,\dd x$ を求めましょう。
\bkpans{$\Bigl[\,xe^x-e^x\,\Bigr]_0^1=(e-e)-(0-1)=1$}

\bkproblemgroup{チャレンジ}{4}

\bkproblem 電流 $i=4t\,e^{-2t}$（A）が流れました。はじめの1秒間に流れた電気の量を求めましょう（$e^{-2}=0.135$）。
\bkpans{$\Bigl[\,-2te^{-2t}\,\Bigr]_0^1+\displaystyle\int_0^12e^{-2t}\,\dd t=-2e^{-2}+(1-e^{-2})=1-3e^{-2}\approx\qty{0.59}{C}$}

\bkproblem $\displaystyle\int_0^1x^2e^{-x}\,\dd x$ を求めましょう（$e^{-1}=0.368$）。
\bkpans{$\displaystyle\int x^2e^{-x}\,\dd x=-e^{-x}(x^2+2x+2)+C$。$\Bigl[\,\cdots\,\Bigr]_0^1=2-5e^{-1}\approx0.160$}

\bkproblem $\displaystyle\int_0^{\pi}\sin^2t\,\cos t\,\dd t$ と $\displaystyle\int_0^{\frac{\pi}{2}}\sin^2t\,\cos t\,\dd t$ を求め、ちがいのわけを考えましょう。
\bkpans{$u=\sin t$。0から $\pi$ では $u:0\to0$ で0。0から $\frac{\pi}{2}$ では $u:0\to1$ で $\dfrac{1}{3}$。$\frac{\pi}{2}$ から $\pi$ では $\cos t<0$ で、面積が打ち消しあう。}

\bkproblem $\displaystyle\int_0^1\frac{x}{x^2+1}\,\dd x$ と $\displaystyle\int_0^{\frac{\pi}{2}}x\sin x\,\dd x$ を求めましょう（$\ln2=0.693$）。
\bkpans{$u=x^2+1$（$1\to2$）で $\dfrac{1}{2}\Bigl[\,\ln u\,\Bigr]_1^2=\dfrac{1}{2}\ln2\approx0.347$。$\Bigl[\,-x\cos x+\sin x\,\Bigr]_0^{\frac{\pi}{2}}=1$}

\bkproblem $\displaystyle\int_{-2}^{2}\sqrt{4-x^2}\,\dd x$ を、$x=2\sin\theta$ とおいて求めましょう。答えが、半径2の半円の面積と同じになることを確かめましょう。
\bkpans{$\displaystyle\int_{-\frac{\pi}{2}}^{\frac{\pi}{2}}4\cos^2\theta\,\dd\theta=4\times\dfrac{\pi}{2}=2\pi$。半円 $\dfrac{1}{2}\times\pi\times2^2=2\pi$}

\begin{furikaeri}
\begin{checklist}
\item 置換積分で、$\dd u$ への書きかえができる。\hfill{\footnotesize\sffamily\color{bkTheme}→\pageref{bk:5-16:chikan}ページ}
\item 定積分の置換で、はんいを書きかえられる。\hfill{\footnotesize\sffamily\color{bkTheme}→\pageref{bk:5-16:hanni}ページ}
\item 部分積分で、$f$ と $g'$ を選んで計算できる。\hfill{\footnotesize\sffamily\color{bkTheme}→\pageref{bk:5-16:bubun}ページ}
\end{checklist}
\end{furikaeri}

\bkfukidashi{delta}{積分の道具がそろった。次は、テイラー展開だ。}
